\documentclass[11pt]{article}
\usepackage{amsmath,amssymb,amsthm}
\usepackage[margin=1.15in]{geometry}
\usepackage{hyperref}
\newtheorem{lemma}{Lemma}
\newtheorem{proposition}{Proposition}
\newtheorem{remark}{Remark}
\newcommand{\C}{\mathbb{C}}
\title{A quotient-Jacobian identity for the full $(1,-1,-2)$-equivariant class\\ containing the July 2026 counterexample to the Jacobian conjecture\\[2mm]
\large Two propositions circulated for expert comment: is this standard?}
\author{Anonymous\thanks{Pseudonymous deposit.
Correspondence: \texttt{minustsquared@gmail.com}. A SHA-256 commitment to
the authors' identity is published in \texttt{COMMITMENT.txt} in the
accompanying deposit.}}
\date{July 2026}
\begin{document}
\maketitle

\noindent\textbf{Purpose.} Let $\C^{*}$ act on $\mathbb{A}^{3}$ with
\emph{coordinate} weights $(1,-1,-2)$, i.e.\
$t\cdot(x,y,z)=(tx,t^{-1}y,t^{-2}z)$. We consider equivariant polynomial
self-maps whose output components, in order, have weights $(-2,-1,1)$; the
counterexample to the Jacobian conjecture announced by Alp\"oge on July 20,
2026 \cite{alpoge2026}, verified in \cite{ulam2026}, is of this type. We
record an explicit quotient calculus for \emph{every} polynomial map with
this equivariance. The conceptual content fits in one sentence: such a map
multiplies the orbit coordinate by a factor $T$, and the two target
invariants involve the first and second powers of the orbit coordinate, so
the induced map on the quotient plane acquires a Jacobian factor of exactly
$T^{2}$. The machinery is standard invariant theory; we record a class-wide
quotient identity that appears not to be stated in the currently circulating
factorization and cubic-model analyses (the precise, machine-verified
relationship to those analyses is given under ``Related formulations''
below). We would be grateful to learn whether either statement
is known, and if so where. The construction of the counterexample itself is
due to \cite{alpoge2026} and, independently, to the derivation published in
\cite{lou2026}; nothing in this note claims otherwise.

\medskip
\noindent\textbf{Setup.} The invariant ring is $\C[x,y,z]^{\C^*}=\C[u,v]$ with
$u=xy$, $v=x^2z$. The weight-$w$ isotypic components relevant here are the
$\C[u,v]$-modules
$A_{-2}=z\,\C[u,v]+y^2\,\C[u,v]$, $A_{-1}=y\,\C[u,v]+xz\,\C[u,v]$,
$A_{1}=x\,\C[u,v]$, with relations $v\,y^2=u^2 z$ and $v\,y=u\,(xz)$.

\begin{lemma}[Normal form]
Every polynomial map $F:\C^3\to\C^3$ equivariant with component weights $(-2,-1,1)$
is of the form
\[
F_1=zP(u,v)+y^2Q(u,v),\qquad F_2=yR(u,v)+xzS(u,v),\qquad F_3=xT(u,v),
\]
the representation being unique up to the gauge
$(P,Q)\mapsto(P+u^2H_1,\,Q-vH_1)$, $(R,S)\mapsto(R+vH_2,\,S-uH_2)$, and unique
if one requires $Q\in\C[u]$ and $S\in\C[v]$. If $F$ is a Keller map, then $T\not\equiv0$, the
linear part of $F$ is the invertible anti-diagonal $(\alpha z,\beta y,\gamma x)$
with $-\alpha\beta\gamma=\det DF$, and $\alpha=\beta=\gamma=1$ may be arranged by
diagonal linear changes of source and target. Precisely, the equivalence used is
left--right equivalence $F\mapsto D_{\mathrm{tgt}}\circ F\circ D_{\mathrm{src}}$
with $D_{\mathrm{src}},D_{\mathrm{tgt}}$ invertible diagonal (such maps commute
with the torus action, hence preserve the equivariance class); the numerical
value of the constant $\det DF$ is \emph{not} preserved --- it rescales by the
product of the six diagonal entries --- but the Keller property is, and no
determinant-one constraint is imposed. The known counterexample has normal form
$P=(1+u)^3$, $Q=(1+u)(4+3u)$, $R=1+12u+9u^2+(6+3u)v$, $S=3$, $T=2-3u-v$.
\end{lemma}

\begin{proof}
A monomial $x^iy^jz^k$ has weight $i-j-2k$. For weight $1$: $i=j+2k+1$, so the
monomial is $x\,u^jv^k$. For weight $-1$: if $j\ge1$ it is $y\,u^{j-1}v^k$; if
$j=0$ then $k\ge1$ and it is $xz\,u^{j}v^{k-1}$; both representations exist iff
$j,k\ge1$, whence the single relation $vy=u(xz)$, and the gauge as stated. In
module terms: the kernel of $\C[u,v]^{2}\to A_{-1}$, $(R,S)\mapsto yR+xzS$, is
the free rank-one module generated by $(v,-u)$, which is exactly the gauge
$(R,S)\mapsto(R+vH_2,\,S-uH_2)$. Weight
$-2$ is identical with generators $z,y^2$ and relation $vy^2=u^2z$; there the
kernel of $(P,Q)\mapsto zP+y^{2}Q$ is generated by $(u^{2},-v)$. The gauge
shifts $Q$ by $-vH_1$ and $S$ by $-uH_2$, so the $v$-multiples of $Q$ and the
$u$-multiples of $S$ can be removed, arranging $Q\in\C[u]$ and $S\in\C[v]$. This
normal form is unique: every monomial of $y^2Q$ with $Q\in\C[u]$ is free of $z$
while every monomial of $z\,\C[u,v]$ contains $z$, and every monomial of $xzS$
with $S\in\C[v]$ is free of $y$ while every monomial of $y\,\C[u,v]$ contains
$y$, so no further cancellation is possible and the residual gauge is $0$. For the Keller statements: $\det DF$ is constant, hence
equals its value at $0$, the determinant of the linear part; the only equivariant
linear monomials are $z,y,x$ in the three slots; and $T\equiv0$ would force
$F_3\equiv0$.
\end{proof}

\begin{proposition}[Square-Jacobian identity]
Let $F$ be as above and let $G:\ \C^2\to\C^2$ be the induced map on the
categorical quotient $\mathbb A^3/\!/\C^*\cong\mathbb A^2$, explicitly
\[
G(u,v)=\bigl(T\,(uR+vS),\ T^2\,(vP+u^2Q)\bigr),
\]
which is the expression of the target invariants $F_2F_3$ and $F_1F_3^2$ in the
source invariants. Then, identically in $\C[u,v]$,
\[
\det DG \;=\; -\,T^{2}\,\det DF .
\]
In particular $F$ is Keller with $\det DF=c$ iff\/ $\det DG=-c\,T^2$: the quotient
map of an \'etale equivariant map has Jacobian a nonzero constant times a perfect
square. Its critical locus is the scheme $V(T^{2})$, with reduced support $V(T)$;
when $T$ is nonconstant this is a divisor which $G$ contracts to the origin
(when $T$ is a nonzero constant the critical locus is empty), and the
multiplicity two is forced by the mechanism described in the opening paragraph.
\end{proposition}

\begin{proof}
On $x\neq0$, $(x,u,v)$ are coordinates and
$dx\wedge du\wedge dv=x^{3}\,dx\wedge dy\wedge dz$, since
$du\wedge dv=(x\,dy+y\,dx)\wedge(2xz\,dx+x^{2}dz)$ contributes only
$x^{3}dy\wedge dz$ after wedging with $dx$. On the target, with coordinates
$(p,q,r)$ and invariants $u'=qr$, $v'=pr^{2}$, the same computation gives
$du'\wedge dv'\wedge dr=-r^{3}\,dp\wedge dq\wedge dr$. Pulling back by $F$:
the left side is
$d(G_1\!\circ\!\pi)\wedge d(G_2\!\circ\!\pi)\wedge d(xT)
=(\det DG)\,du\wedge dv\wedge(T\,dx+x\,dT)
=(\det DG)\,T\,du\wedge dv\wedge dx$,
because $dT\in\C[u,v]\,du+\C[u,v]\,dv$; the right side is
$-(xT)^{3}(\det DF)\,dx\wedge dy\wedge dz$. Comparing via
$du\wedge dv\wedge dx=dx\wedge du\wedge dv=x^{3}dx\wedge dy\wedge dz$
yields the polynomial identity $(\det DG)\,T\,x^{3}=-x^{3}T^{3}\det DF$ on the
dense open set $x\neq0$, hence in $\C[x,u,v]$. If $T\equiv0$ then $G=(0,0)$
and both sides of the asserted identity vanish. If $T\not\equiv0$, cancel
$x^{3}T$ in the integral domain $\C[x,u,v]$ to obtain
$\det DG=-T^{2}\det DF$ in $\C[u,v]$.
\end{proof}

\begin{remark}
Conceptually this is the transformation law of the quotient canonical form: contracting
$dx\wedge dy\wedge dz$ with the Euler field $E=x\partial_x-y\partial_y-2z\partial_z$
of the action produces the horizontal $2$-form descending to $du\wedge dv$ (up to the
orbit coordinate), and $T$ is the multiplier by which $F$ scales the orbit direction,
entering squared through the two target invariants of weights involving $F_3$ and
$F_3^2$. We regard the abstract principle as standard; our question concerns the
explicit formula and its use below.
\end{remark}

\begin{proposition}[Generic degree]
If $F$ is dominant and $T\not\equiv0$ (both automatic for Keller maps), then
\[
[\C(x,y,z):\C(F_1,F_2,F_3)]=[\C(u,v):\C(G_1,G_2)],
\]
i.e.\ $\mu(F)=\mu(G)$. Here $[\,\cdot:\cdot\,]$ denotes field-extension
degree; over $\C$, for a dominant map, this coincides with the cardinality of
the generic fiber (and with the separable degree, characteristic zero making
all extensions separable), so no ambiguity among the usual notions of
``generic degree'' arises.
\end{proposition}

\begin{proof}
Set $K=\C(u,v)$ and $L=\C(x,y,z)$; since $y=u/x$ and $z=v/x^2$, $L=K(x)$. Let
$w=F_3=xT(u,v)$. Since $x$ is transcendental over $K$ and $T\in K^{\times}$,
$w=xT$ is transcendental over $K$ and $x=w/T\in K(w)$, so $L=K(x)=K(w)$ is a
purely transcendental extension of $K$. From $G_1=F_2F_3$
and $G_2=F_1F_3^{2}$ we get $F_2=G_1/w$ and $F_1=G_2/w^{2}$, hence
$\C(F_1,F_2,F_3)=\C(G_1,G_2,w)$. Dominance of $F$ gives
$\operatorname{trdeg}_{\C}\C(F_1,F_2,F_3)=3$; as
$\C(F_1,F_2,F_3)=\C(G_1,G_2,w)$ and $w$ contributes at most one to the
transcendence degree, $\operatorname{trdeg}_{\C}\C(G_1,G_2)=2$, i.e.\ $G$ is
dominant, $w$ is transcendental over $\C(G_1,G_2)$, and both displayed
extension degrees are finite. Adjoining the transcendental $w$ preserves
relative degree, so
$[L:\C(F_1,F_2,F_3)]=[K(w):\C(G_1,G_2)(w)]=[K:\C(G_1,G_2)]$.
\end{proof}

\begin{remark}
For the known counterexample, $G$ has Jacobian $2(2-3u-v)^{2}$ and generic
degree $3$, consistent with the cubic fiber description of
\cite{ulam2026}. The degree claim is certified as follows (script
\texttt{check\_quotient\_degree.py} and transcript
\texttt{quotient\_degree\_check.log} in the accompanying materials): the lex
Gr\"obner basis of $\langle G_1-s,\ G_2-t\rangle$ over the exact function
field $\mathbb{Q}(s,t)$, variables $u>v$, is triangular, consisting of a
degree-$3$ eliminant in $v$ and one element linear in $u$; the eliminant is
squarefree (its gcd with its $v$-derivative is constant), so the generic
fiber has exactly $3$ distinct points and
$[\C(u,v):\C(G_1,G_2)]=3$. Fiber-by-fiber
statements over $T=0$ and on $x=0$, where the action degenerates, require separate
treatment and are not claimed here.
\end{remark}

\medskip
\noindent\textbf{Related formulations.} A perfect-square Jacobian factor
appears in two independent, essentially contemporaneous analyses of the
specific counterexample, and the relationship of Proposition~1 to both can
now be stated exactly rather than conjecturally. The derivation summary
published by Lou \cite{lou2026} (produced by an OpenAI Codex model that
independently rediscovered the map) factors a cubic as $L\cdot Q$ with
$L=xt+\beta$, $Q=\gamma t^{2}+\delta t+\varepsilon$, obtains the
five-dimensional coefficient-map identity
$\det\partial(c_3,c_2,c_1,c_0,\rho)/\partial(x,\beta,\gamma,\delta,\varepsilon)
=-\rho^{2}$ (its Lemma~1), and passes to the slice $c_2=\rho=1$ via an
explicit polynomial chart $(x,y,z)$ (its Proposition~2). That chart is
precisely the normal-form coordinate system of our Lemma~1, and the following
dictionary is machine-verified in exact arithmetic (script
\texttt{check\_lou\_dictionary.py}, transcript
\texttt{lou\_dictionary\_check.log}, in the accompanying materials):
$\beta=1+u$, and identically $T=2\gamma$, so the contracted divisor $V(T)$ of
Proposition~1 is $\{\gamma=0\}$, the locus where the quadratic factor drops
degree --- the nonproperness divisor. In particular $T$ does \emph{not}
correspond to $\rho$: on the slice $\rho\equiv1$ identically, and $\rho$ ---
like the ramification divisor of the symmetric-product construction obtained
by Jiang working with ChatGPT and recorded by Sawin \cite{sawin2026}, both
measuring coincidence of factors --- is spent by the normalization, while $T$
survives to the quotient. The two squares are therefore \emph{distinct
objects in one geometry}: $-\rho^{2}$ upstream (finite ramification, removed
by the resultant slice) and $-T^{2}$ downstream (the orbit-direction
multiplier squared, contracted by $G$). The binary-cubic fiber calculus of
\cite{ulam2026} is the per-map instance of the resulting degree-$3$ quotient
geometry.

\medskip
\noindent\textbf{Questions for the reader.}
(1) Are the normal form and the identity of Proposition~1 recorded anywhere, in
this or equivalent notation? (2) Does the reading of the counterexample as a
quotient-plane map with perfect-square Jacobian contracting $T=0$ appear in any
current treatment? (3) An earlier draft posed here the question of whether
$T$ matches Lou's $\rho$; it is resolved negatively by the verified
dictionary above ($T=2\gamma$, while $\rho\equiv1$ on the slice). Is the
resulting two-square pattern --- a coincidence square spent by normalization,
an infinity square contracted on the quotient --- a known phenomenon for
torus quotients of \'etale maps? (4) We are also investigating low-degree classification
questions for this equivariant class; a complete, machine-certified
classification through polynomial degree four, with all certificates, solver
transcripts, and reproduction materials, accompanies this note in the same
deposit (see the bundle's \texttt{README} and \texttt{HANDOFF} files for the
precise statements, software versions, and what ``exhaustive'' means there).
Comments on whether the higher-degree program is worth pursuing would be
greatly valued.

\medskip
\noindent\textbf{Acknowledgment.} The quotient formulation was found, and
the exploratory symbolic computations carried out, with the assistance of
Claude Fable (Anthropic). Every mathematical claim above is certified by
machine-checkable artifacts in exact arithmetic --- certificates, solver
transcripts, and verification scripts in the accompanying deposit ---
independent of any AI system's reasoning and reproducible from the supplied
materials.

\begin{thebibliography}{9}
\bibitem{alpoge2026} L.~Alp\"oge, announcement of a counterexample to the
Jacobian conjecture in dimension three, X post, July 20, 2026,
\url{https://x.com/__alpoge__/status/2079028340955197566}.
\bibitem{ulam2026} A counterexample to the Jacobian conjecture, verification
paper, July 20, 2026, \url{https://www.ulam.ai/research/jacobian.pdf}.
\bibitem{lou2026} A.~Lou, \emph{Deriving an explicit polynomial
counterexample to the Jacobian conjecture}, derivation summary produced by an
internal OpenAI Codex model, July 20, 2026,
\url{https://aaronlou.com/jacobian_counterexample_derivation.pdf};
announcement: \url{https://x.com/aaron_lou/status/2079218392452530249}.
\bibitem{sawin2026} W.~Sawin, comment (July 20, 2026, 12:19pm) on D.~Speyer,
\emph{The new counterexample to the Jacobian conjecture}, Secret Blogging
Seminar, recording a symmetric-product construction obtained by A.~Jiang
working with ChatGPT,
\url{https://sbseminar.wordpress.com/2026/07/20/the-new-counterexample-to-the-jacobian-conjecture/\#comment-27948}.
\end{thebibliography}
\end{document}
